% !TeX root = review-main.tex

\begin{center}
{\Large\textbf{第3回 復習テスト}}
\end{center}

\vspace{0.2em}
\begin{flushright}
氏名：\underline{\hspace{14em}}
\end{flushright}

\vspace{0.3em}
\noindent 以下の問題は、配布テキストに沿って解答すること。(10 $\times$ 10点)

\vspace{0.5em}

\noindent\textbf{【問1】} 次の英文を疑問文にしなさい。

「He can play the guitar very well.」

\testblank[Can he play the guitar very well?]{0.92\linewidth}
\testqsep

\noindent\textbf{【問2】} 問1の英文を\uwave{9語}で書き換えなさい。

\testblank[He is able to play the guitar very well.]{0.92\linewidth}
\testqsep

\noindent\textbf{【問3】} 次の英文を否定文にしなさい。

「She will visit Kyoto next week.」

\testblank[She won't visit Kyoto next week.]{0.92\linewidth}
\testqsep

\noindent\textbf{【問4】} 問3の英文を\uwave{9語}で書き換えなさい。

\testblank[She is not going to visit Kyoto next week.]{0.92\linewidth}
\testqsep

\noindent\textbf{【問5】} 次の日本語を主語と動詞を含む英文にしなさい。

「あなたのペンを借りてもいいですか。」

\testblank[May I use your pen?]{0.92\linewidth}
\testqsep

\noindent\textbf{【問6】} 次の日本語を主語と動詞を含む英文にしなさい。

「彼は疲れているにちがいない。」

\testblank[He must be tired.]{0.92\linewidth}
\testqsep

\noindent\textbf{【問7】} 次の英文を和訳しなさい。

「You must not be late for school.」

\testblank[学校に遅れてはいけないよ。]{0.92\linewidth}
\testqsep

\noindent\textbf{【問8】} 問7の英文を\uwave{5語}で書き換えなさい。

\testblank[Don't be late for school.]{0.92\linewidth}
\testqsep

\noindent\textbf{【問9】} 次の英文を和訳しなさい。

「I would like to visit Hokkaido.」

\testblank[北海道を訪れたいです。]{0.92\linewidth}
\testqsep

\noindent\textbf{【問10】} 自分が将来、できるであろうことについて、主語と動詞を含む英文1文を作りなさい。

\testblank[I will be able to swim.]{0.92\linewidth}
\testqsep

\begin{flushright}
\textbf{得点} \underline{\hspace{8em}} ／ 100点
\end{flushright}
